% Section content template for individual Educative lessons
% This file contains the actual content for a specific section
% Generated from Educative API section content

% Section: Quiz: Object-Oriented Design
% Section ID: 4504644179197952
% Section Slug: quiz-object-oriented-design
% Generated: 2025-12-08T22:23:32.196723

\subsection{Arrange the sequence diagram messages}\label{7yHmnfPaWNT9wtKvdAXxa}
\\
\textbf{Scenario:} You are a student wishing to create an account on your university's web portal using some credentials. After signing up, you also want to make sure that your account is verified.
\\
For the given scenario, arrange the following messages in order to appear in a sequence diagram. The messages are given in the following format:

\[
<<ObjectName>>\space:\space <<message>>
\]

\begin{quote}
\textbf{Note:} The object name here refers to the object that sends the message. Here, we assume that the student enters the correct credentials.
\end{quote}


\subsection{MCQs}\label{SieVIhYB2lT88NY28ePjZ}
\\
Challenge yourself by solving the following quiz questions.

\subsection*{Object-Oriented Design}
\vspace{0.3cm}
\noindent
\textbf{Question 1:} Which statement about the given class diagram is \emph{true}?
\\
\includegraphics[width=4.16667in,height=\textheight,keepaspectratio]{/api/collection/10370001/5583710957338624/page/4504644179197952/image/4945256430960640?page_type=collection_lesson}
\\[0.2cm]
\begin{itemize}
 \item An instance of class \texttt{A} can be associated with multiple instances of class \texttt{D}, and each instance of class \texttt{D} can be associated with exactly two instances of class \texttt{A}.
 \item An instance of class \texttt{A} can be associated with multiple instances of class \texttt{D}, and each instance of class \texttt{D} can be associated with any number of instances of class \texttt{A}.
 \item \textbf{[CORRECT]} An instance of class \texttt{A} can be associated with exactly two instances of class \texttt{D}, and each instance of class \texttt{D} can be associated with any number of instances of class \texttt{A}.
\\[0.1cm]
\textit{Explanation:} 
Class \texttt{A} has a two-way association with class \texttt{D}, and it has child classes \texttt{B} and \texttt{C}.
 \item An instance of class \texttt{A} can be associated with exactly two instances of class \texttt{D}, and each instance of class \texttt{D} can be associated with exactly one instance of class \texttt{A}.
\\[0.1cm]
\textit{Explanation:} 
One object of {{\(D\)}} is associated with multiple objects of {{\(A\)}}.
\end{itemize}
\vspace{0.3cm}
\noindent
\textbf{Question 2:} Which statement about class diagrams is \emph{true}?
\\[0.2cm]
\begin{itemize}
 \item Every class in a class diagram must have both attributes and operations.
 \item Classes always have a section showing a textual description.
 \item \textbf{[CORRECT]} Operations may have parameters and return values.
 \item All attributes in a class must have the same data type.
\end{itemize}
\vspace{0.3cm}
\noindent
\textbf{Question 3:} \textbf{(Select all that apply.)} Which statement(s) about actors in a use case diagram is \emph{true}?
\\
\textit{(Select all that apply.)}
\\[0.2cm]
\begin{itemize}
 \item \textbf{[CORRECT]} Actors can be linked to each other by inheritance.
 \item Actors communicate with other actors.
 \item \textbf{[CORRECT]} Actors can be linked to abstract and non-abstract use cases via associations.
 \item Actors interact with the system in the form of \textless\textless{}\emph{include}\textgreater\textgreater{} relationships.
\end{itemize}
\vspace{0.3cm}
\noindent
\textbf{Question 4:} Which statement about the following use case diagram is \emph{true}?
\\
\includegraphics[width=2.60417in,height=\textheight,keepaspectratio]{/api/collection/10370001/5583710957338624/page/4504644179197952/image/6322483638829056?page_type=collection_lesson}
\\[0.2cm]
\begin{itemize}
 \item Use case {{\(B\)}} can extend use case {{\(A\)}}.
 \item Use case {{\(A\)}} always has to invoke use case {{\(B\)}}.
 \item Use case {{\(A\)}} cannot be executed without use case {{\(B\)}}.
 \item \textbf{[CORRECT]} Use case {{\(B\)}} may or may not invoke use case {{\(A\)}}.
\end{itemize}
\vspace{0.3cm}
\noindent
\textbf{Question 5:} Consider the following activity diagram:
\\
\includegraphics[width=2.08333in,height=\textheight,keepaspectratio]{/api/collection/10370001/5583710957338624/page/4504644179197952/image/6706295438835712?page_type=collection_lesson}
\\
What is \emph{not} the correct order of action sequences that are \emph{not} possible during the execution of the activity?
\\[0.2cm]
\begin{itemize}
 \item A1 → A2
 \item A1 → A3 → A4
 \item A1 → A5 → A6
 \item \textbf{[CORRECT]} A1 → A3 → A6 → A5 → A4
\\[0.1cm]
\textit{Explanation:} 
We can't go from A3 to A6, A6 to A5 and A5 to A4.
\end{itemize}
\vspace{0.3cm}
\noindent
\textbf{Question 6:} In an activity diagram, which component combines multiple non-concurrent flows?
\\[0.2cm]
\begin{itemize}
 \item \textbf{[CORRECT]} Merge
 \item Join
 \item Decision
 \item Fork
\end{itemize}

% End of section content